\chapter{The next research questions}
\label{ch:research-program}

Run the pieces on one real document and we can answer a small question: do they
work together? We still will not know whether the workflow reduces review burden
across genres, writers, and readers. To answer that larger question, the team
must first measure the objects correctly, then watch people use the workflow,
and only then estimate its effect on outcomes.

A demonstration earns the right to ask the larger question. It does not answer
it in advance.

The four rungs in Figure~\ref{fig:research-ladder} make that order explicit.
An attractive reader result cannot rescue a source-correspondence failure
below it.

\hvFigureResearchLadder

\section{Three claims that must not be collapsed}

The team will make three increasingly strong claims. The first is an
\emph{availability claim}: a parser, diagnostic, or review record can be
installed and run on the target document formats. The second is an
\emph{operational claim}: the component preserves the objects and locations
needed by an author, and its output does not impose more review work than it
removes. The third is an \emph{effectiveness claim}: the complete workflow
changes a reader-facing outcome relative to the incumbent bundle.

Each claim needs a different observation. A package README can establish
availability. A fixture benchmark and an observed review session can establish
operational suitability. Only a comparative study with a declared comparator,
reader rubric, and uncertainty can support effectiveness. The same hierarchy
applies to the language-model adapter: plausible questions show availability;
calibrated abstentions and errors show operational suitability; neither proves
that the adapter improves a document.

\begin{table}[H]
\centering
\small
\begin{tabularx}{\textwidth}{@{}p{0.19\textwidth}L p{0.25\textwidth}@{}}
\toprule
Claim & Minimum evidence & Decision it permits \\
\midrule
Availability & Installation record, version, license, supported input, and a
smoke run. & Keep the component in the candidate inventory. \\
Operational suitability & Held-out fixtures, protected-object preservation,
runtime, failure modes, privacy review, and author burden. & Use the component
in the feasibility workflow. \\
Workflow effectiveness & Preregistered comparison against the incumbent,
human reader outcomes, burden, meaning checks, and uncertainty. & Consider
adoption for the measured population and task class. \\
Economic value & Observed volume, loaded rates, avoided rounds, error costs,
and sensitivity. & Fund expansion, revise the offer, or stop. \\
\bottomrule
\end{tabularx}
\caption{Evidence strength rises with the decision at stake.}
\label{tab:research-claim-ladder}
\end{table}

An inventory earns space in the main argument when it explains a choice or a
risk. A long list cannot stand in for a benchmark, so the full package list
stays in an appendix while the main chapters follow the decisions it informs.

\section{The measurement ladder}

The first research stage is a ladder of observations, not a feature roadmap.
Each rung names the object the team will measure and the failure that would
stop the next step.

\subsection{Rung 1: source correspondence}

For every target format, the system must map rendered passages back to source
locations and identify protected objects. The fixtures should include ordinary
prose, equations with macros, labels, citations, tables, quotations, code,
comments, malformed input, and deliberately duplicated identifiers. The test
does not ask whether a parser is elegant. It asks whether an author can locate
the finding and whether a before-and-after comparison can tell what changed.

Let $O$ be the set of protected objects in a source version and let
$\pi(o)$ be the rendered location associated with object $o$. A parser pass is
usable only if it preserves a partial map
\[
  \pi: O \longrightarrow \{\text{source spans}\}\times
  \{\text{rendered anchors}\}
\]
for every object that the product promises to protect. Missing maps are not
silently imputed. They are an explicit ``unparsed'' result that blocks the
protected comparison.

\subsection{Rung 2: diagnostic validity}

The second stage asks whether a finding corresponds to an observable property
of the current document. Register leakage can be counted against a declared
lexicon and context rule. Repeated structure can be measured against a
document-level baseline. Citation identity can be checked against metadata.
These are appropriate places for precision, recall, and abstention measures.

The stage does not turn every editorial judgment into a classification task.
Whether a qualification belongs in a paragraph, whether a term is familiar to
the named reader, and whether a conclusion is proportionate to its evidence
are questions for an author or reader. The diagnostic may attach a question
and its evidence; it must not pretend that a threshold has answered it.

\subsection{Rung 3: workflow burden}

A finding has a cost. The author must read it, decide whether it applies, and
possibly revise the source. The reader may later encounter the same passage
and request a new change. The feasibility run should record the time spent in
each activity, the number of findings opened, the number accepted or rejected,
the number of rollbacks, and the reasons for abstention.

If $q$ is the number of findings shown, $a$ the number accepted, and $t_j$ the
time spent on finding $j$, a simple burden measure is
\[
  B=\sum_{j=1}^{q}t_j + t_{\mathrm{recheck}}.
\]
The expression is not a universal utility function. It is a bookkeeping
device that prevents a high precision score from hiding a queue that nobody
can afford to inspect. Report $B$ alongside the findings that led to a reader
revision and the findings that were rejected as irrelevant.

\subsection{Rung 4: reader reconstruction}

Give the rendered document to a reader without the project history and ask five
questions: what does it answer, which objects carry the answer, what evidence
supports it, what remains uncertain, and what action follows? A reader can
understand a paragraph without agreeing with its conclusion, so the rubric
separates comprehension, credibility, and willingness to act.

The first reader exercise is timed and conducted without project history. The
reader receives
the rendered document, the declared role, and the decision question, but not
the source map, findings, or private project history. A second exercise may
show the review packet to measure whether the packet makes revision easier.
Keeping the exercises separate distinguishes the product's author interface
from the outcome it is meant to improve.

\section{The corpus is a scientific instrument}

Do not build the corpus from convenient examples alone. It must represent the
documents on which the product will be used and preserve why a finding was
accepted or rejected. The existing cases supply valuable repair pairs, but they
are selected cases. The evaluation corpus should add prospective samples
from at least three genres: an investment or funding proposal, a technical
research note, and a long-form methodological or monograph chapter.

Each unit should carry a document-level reading brief:

\begin{enumerate}
\item the writer's baseline experience and role;
\item the intended reader and the decision the document supports;
\item the source format, length, equation and citation counts, and sensitivity
      class;
\item the incumbent workflow, including which tools and human review were
      actually used;
\item protected-object annotations and known substantive constraints; and
\item named exemplar documents, known reader vocabulary, and declared
      exemption zones for first-use checks;
\item pacing profile, first-use findings, author dispositions, and the
      detector version that produced them; and
\item the reader outcome, revision history, and adjudicated failure reasons.
\end{enumerate}

The split must be made at the document lineage, not at the paragraph level.
Otherwise a repeated passage can leak from tuning into evaluation and produce
an implausibly clean result. A held-out document should remain held out from
lexicon selection, prompt design, model calibration, and author training.

\paragraph{A corpus split that can fail honestly.}
Take one long technical chapter and its three revised versions. Keeping one
version in training and another in test is not a held-out evaluation: the
reading brief, equations, and private terminology have already leaked. The
lineage belongs to one partition. If the held-out set becomes too small, the
correct response is to collect more documents, not to split the same chapter
more finely.

The labels need two levels. A deterministic label records whether a protected
object was preserved, whether a citation identity resolved, or whether a
declared pattern occurred. An adjudicated label records whether the finding
was useful, whether the revision changed the intended meaning, and whether the
reader could act. The first level supports repeatable software tests; the
second supports product claims.

The HPO case adds a practical unit of work. Select one chapter, run the
deterministic meters, let an author repair or dismiss the ranked findings,
rebuild the chapter, and then carry the updated brief into the next unit. A
whole-document pass before the unit is understood hides which repair caused a
reader response and makes a false positive expensive to remove. The corpus
therefore stores unit ids and brief versions alongside document ids.

\section{The comparative program}

The comparative study should begin only after the feasibility run has fixed
the workflow and the comparator. The incumbent is not ``no tool.'' It is the
combination the author would otherwise use: editor, general language model if
permitted, version control, citation checks, and manual review. Removing those
components to make Humanvoice look better would answer the wrong question.

For each document task, randomize or counterbalance the order in which the
incumbent and Humanvoice versions are prepared and read. Blind the reader to
condition where the document and disclosure rules permit. Record protocol
deviations rather than silently excluding difficult tasks. The primary outcome
is the difference in feedback rounds to acceptance; secondary outcomes include
reader time, author burden, protected-object incidents, requested substantive
changes, and confidence in the decision.

Heterogeneity is part of the question. Report results by writer experience,
genre, reader role, document length, equation density, and whether a model
adapter was used. If an average improvement comes entirely from short business
documents while long technical chapters become harder to review, the average
does not justify expansion to the latter population.

\subsection{Calibration and model judges}

Model judges may reduce the cost of screening a large set of pairwise
comparisons, but they are not the acceptance authority. Before use, calibrate
them on human-labelled examples with randomized order, length controls, and
multiple model families where practical. Measure position bias, verbosity
bias, self-preference, and sensitivity to protected mathematical content.
Require abstention when the judge cannot identify the reading brief or when
the two versions differ in a protected object.

The judge's output should be treated as a noisy measurement. If $Y$ is the
human reader label and $J$ the model label, report the confusion matrix and
calibration curve rather than only an agreement percentage. A high agreement
on easy surface edits can coexist with systematic failure on claim scope or
technical uncertainty. The human-labelled subset must therefore oversample
the cases most likely to expose those failures.

\section{From feasibility to an operating service}

If the comparative result is favorable, adoption still requires an operating
model. A document owner must know where source text is processed, how long
snapshots and rendered pages are retained, who can see the reader decision,
and how a failed parse is recovered. The product should expose a local-first
path and an explicit remote-model boundary. Disclosure is part of the record:
the author can state whether a model suggested a question, whether a human
accepted the revision, and whether any confidential excerpt left the local
environment.

The service should also define what it does when the product is unavailable.
An author must be able to continue with the incumbent workflow, preserve the
source version, and import the eventual decision later. A dependency that
blocks ordinary writing is a deployment risk independent of model quality.

\begin{table}[H]
\centering
\small
\begin{tabularx}{\textwidth}{@{}p{0.24\textwidth}L L@{}}
\toprule
Operating question & Minimum answer before expansion & Evidence owner \\
\midrule
Confidentiality & Which text, metadata, and rendered pages may leave the local
boundary? & Project owner and security reviewer \\
Retention & Which snapshots, findings, and reader responses are retained, and
for how long? & Product owner \\
Continuity & Can the author complete the document if Humanvoice or a model
adapter is unavailable? & Product engineer \\
Disclosure & What assistance is disclosed to the reader, funder, or publisher?
& Author and institutional policy owner \\
Escalation & Who decides when a protected comparison blocks release or a
reader requests substantive repair? & Named document owner \\
\bottomrule
\end{tabularx}
\caption{Operating questions that remain after a favorable experiment.}
\label{tab:research-operating-model}
\end{table}

\section{Alternatives if the full product does not earn adoption}

The research program should preserve useful partial outcomes. If protected
source correspondence works but the reader effect is null, the repository may
still have a valuable technical comparison library. If deterministic
diagnostics reduce author burden but the full packet does not change reader
acceptance, the project may become an internal editing instrument rather than
a product sold on reader outcomes. If the privacy boundary prevents model
assistance, the deterministic and human-review path can remain the default.

These are not consolation prizes to be declared after a failed study. They are
pre-specified alternatives that keep the interpretation honest. A null result
for the complete workflow does not establish that every component is useless;
it establishes that the bundle did not earn the claimed decision under the
tested conditions. Conversely, a positive result on a narrow corpus does not
support universal deployment. Expansion follows the population for which
the evidence was collected.

\section{The decision after twelve weeks}

At the end of the feasibility horizon, the owner should receive one compact
decision package with the detailed records behind it:

\begin{enumerate}
\item a replayable source-to-reader run on a live document;
\item parser and protected-object fixture results, including failures;
\item the adjudicated corpus description and held-out split;
\item an author-burden and reader-time account;
\item a calibration report for any model judge or editor adapter;
\item an observed cost worksheet with sensitivity to volume and rates; and
\item a recommendation to proceed, revise, or stop, with the evidence that
      would change that recommendation.
\end{enumerate}

The package is deliberately more substantial than a status report. It gives
the sponsor enough detail to challenge an assumption, reproduce a failure,
and choose the next experiment. That is the standard the product itself must
meet: a reader should not have to trust a fluent conclusion when the underlying
objects and decision path can be shown.

The next question is organizational rather than statistical. Suppose the
feasibility packet is reproducible, the protected objects survive, and the
reader rubric can be administered. Who will use the system, what incumbent
work will it replace, and what privacy and pricing conditions make repeated use
reasonable? Those questions are not a promise that the workflow has earned a
market. They are the observations needed to turn a successful experiment into
an accountable operating choice. Chapter~\ref{ch:adoption-scale} answers them
by separating buyer, author, editor, and reader roles, then setting a scale
rule that follows the evidence rather than a projection.
